% Repository-ready section file.
% Theorem environments defined in main.tex preamble.
%
% The status conventions in this file are:
%   - Theorems/Propositions/Lemmas with proofs are formal or representation-theoretic statements
%     proved in this section.
%   - Conjectures are physically motivated or modular-completion statements.
%   - The final ``provisional'' subsection records exact low-mode coefficient identities discussed in
%     the working thread, but labels them explicitly as unverified and not used elsewhere.

\providecommand{\A}{\mathcal A}
\providecommand{\B}{\mathcal B}
\providecommand{\C}{\mathbb C}
\providecommand{\D}{\mathcal D}
\providecommand{\E}{\mathcal E}
\providecommand{\F}{\mathcal F}
\providecommand{\G}{\mathfrak G}
\providecommand{\M}{\mathcal M}
\providecommand{\N}{\mathcal N}
\providecommand{\Q}{\mathbb Q}
\providecommand{\R}{\mathbb R}
\providecommand{\Z}{\mathbb Z}
\providecommand{\mc}{\mathfrak m}
\providecommand{\h}{\hbar}
\providecommand{\id}{\mathrm{id}}
\providecommand{\Hom}{\operatorname{Hom}}
\providecommand{\MC}{\operatorname{MC}}
\providecommand{\Cone}{\operatorname{Cone}}
\providecommand{\Obs}{\operatorname{Obs}}
\providecommand{\Sym}{\operatorname{Sym}}
\providecommand{\Tr}{\operatorname{Tr}}
\providecommand{\Span}{\operatorname{Span}}
\providecommand{\End}{\operatorname{End}}
\providecommand{\Ext}{\operatorname{Ext}}
\providecommand{\Defcyc}{\operatorname{Def}_{\mathrm{cyc}}}
\providecommand{\Ran}{\operatorname{Ran}}
\providecommand{\gr}{\operatorname{gr}}
\providecommand{\Hb}{H_b^{0,\bullet}(S^3)}
\providecommand{\HCR}{H_{\mathrm{CR}}^{0,\bullet}(S^3)}
\providecommand{\delbar}{\bar\partial}
\providecommand{\delbarb}{\bar\partial_b}
\providecommand{\ot}{\widehat\otimes}
\providecommand{\ov}[1]{\overline{#1}}

\section{Modular twisted/celestial holography from first principles}
\label{sec:modular-twisted-celestial-first-principles}

The modular trace principle (Principle~\ref{princ:modular-trace}) extends the disk-level projections of $\Theta^{\mathrm{oc}}$ (boundary chiral algebras, spectral $R$-matrices) to all genera; the genus tower is the all-genus modular completion data.

\subsection{The problem, restated at its true level}
\label{subsec:problem-restated-true-level}

Kaluza--Klein reduction of $6$d holomorphic theories produces $3$d holomorphic--topological theories whose boundary chiral algebras coincide with defect chiral algebras of the unreduced model \cite{ZengTwistedHolography}.
The bulk chiral algebra of the holomorphic twist is commutative with a shifted Poisson bracket, and boundary chiral algebras carry higher $A_\infty$-type operations \cite{CostelloDimofteGaiottoBoundary}.
Regularized Feynman diagrams satisfy quadratic identities controlling holomorphic--topological OPE data at all arities \cite{GaiottoKulpWuHigher}.
Chiral quadratic duality is governed by Maurer--Cartan elements in $A^!\otimes B$, with curvature entering through curved chiral DG structures \cite{GuiLiZengQuadraticDuality}.

These results force a reorganization: the boundary vertex algebra is the genus-$0$ truncation of a full modular Maurer--Cartan lift.
\begin{quote}
\emph{Twisted/celestial holography is a modular lifting problem for the boundary
chiral algebra.}
\end{quote}

\subsection{The abstract modular package}
\label{subsec:abstract-modular-package}

The formal structure of modular completion is as follows.

\begin{definition}[Complete genus-filtered cyclic deformation package]
\label{def:complete-genus-filtered-package}
A \emph{complete genus-filtered cyclic deformation package} consists of the following data.
\begin{enumerate}[label=\textup{(\roman*)}]
\item A complete $\h$-adically filtered graded vector space
\[
\mathfrak g^{\mathrm{mod}}=\prod_{g\ge 0}\h^g\mathfrak g^{(g)}.
\]
\item A family of multilinear operations
\[
\ell_k^{(g)}: \Sym^k(\mathfrak g^{\mathrm{mod}})\longrightarrow
\mathfrak g^{\mathrm{mod}}[2-k],
\qquad g\ge 0,\ k\ge 1,
\]
which together define a complete cyclic $L_\infty$-algebra.
\item A distinguished genus-$0$ Maurer--Cartan element
\[
\pi\in \mathfrak g^{(0),1}
\]
solving the genus-$0$ Maurer--Cartan equation
\[
\sum_{k\ge 1}\frac1{k!}\,\ell_k^{(0)}(\pi^{\otimes k})=0.
\]
\end{enumerate}
A \emph{modular lift} of $\pi$ is a formal series
\[
\Pi(\h)=\pi+\sum_{g\ge 1}\h^g\Theta^{(g)}
\]
satisfying the full Maurer--Cartan equation
\begin{equation}
\label{eq:full-modular-mc}
\sum_{r\ge 0}\h^r\sum_{k\ge 1}\frac1{k!}\,
\ell_k^{(r)}\!\bigl(\Pi(\h)^{\otimes k}\bigr)=0.
\end{equation}
\end{definition}

\begin{remark}
In the applications we have in mind, $\mathfrak g^{\mathrm{mod}}$ is the completed cyclic
boundary deformation complex tensored with a genus package such as
$R\Gamma(\ov{\M}_{\bullet,\bullet},\Q)$, or equivalently the graph/Feynman transform of the
genus-$0$ deformation theory.  Definition
\ref{def:complete-genus-filtered-package} separates the formal deformation theory from the
physical realization.
\end{remark}

\begin{conjecture}[Modular twisted/celestial holography; \ClaimStatusConjectured]
\label{conj:modular-twisted-celestial}
Let $(T,\mathfrak b)$ be a twisted or celestial holographic model whose disk-level boundary
algebra $\A_\partial=\A_\partial(T,\mathfrak b)$ is identified, for a transverse boundary
condition, with the universal defect chiral algebra of the unreduced theory.  Then there exists
a complete genus-filtered cyclic deformation package
\[
\mathfrak g^{\mathrm{mod}}_\partial=\prod_{g\ge 0}\h^g\mathfrak g^{(g)}_\partial
\]
and a distinguished genus-$0$ Maurer--Cartan element
\[
\pi_\partial\in \mathfrak g_\partial^{(0),1}
\]
such that:
\begin{enumerate}[label=\textup{(\roman*)}]
\item the disk-level boundary OPE and transferred $A_\infty$-module structure are encoded by
$\pi_\partial$;
\item the full twisted/celestial holographic dictionary is equivalent to a modular lift
\[
\Pi_\partial(\h)=\pi_\partial+\sum_{g\ge 1}\h^g\Theta_\partial^{(g)}
\]
of $\pi_\partial$;
\item the coefficients $\Theta_\partial^{(g)}$ are the genus-$g$/loop/$1/N$ corrections to the
bulk--boundary dictionary;
\item under the appropriate bulk--boundary Koszul/Verdier comparison, this modular lift is the
same deformation problem as the corresponding bulk modular lift.
\end{enumerate}
Equivalently: the known boundary chiral algebra is the genus-$0$ shadow of a canonical
modular Maurer--Cartan element.
\end{conjecture}

\subsection{Order-by-order obstruction theory}
\label{subsec:order-by-order-obstruction-theory}

Higher-genus corrections are obstruction classes in the $d_\pi$-cohomology of the genus-filtered package.

\begin{definition}[Twisted differential and linearized genus-$0$ complex]
\label{def:twisted-differential}
Given $\pi\in\mathfrak g^{(0),1}$ as in Definition
\ref{def:complete-genus-filtered-package}, define
\begin{equation}
\label{eq:twisted-differential}
d_\pi(x):=\sum_{m\ge 0}\frac1{m!}\,\ell_{m+1}^{(0)}(\pi^{\otimes m},x).
\end{equation}
This is the linearization of the genus-$0$ Maurer--Cartan equation at $\pi$.
\end{definition}

\begin{theorem}[The first modular obstruction; \ClaimStatusProvedHere]
\label{thm:first-modular-obstruction}
Let $\mathfrak g^{\mathrm{mod}}$ be a complete genus-filtered cyclic deformation package and let
$\pi\in \mathfrak g^{(0),1}$ be its genus-$0$ Maurer--Cartan element.  Define the genus-$1$
source term
\begin{equation}
\label{eq:omega-one}
\omega_1(\pi):=\sum_{k\ge 1}\frac1{k!}\,\ell_k^{(1)}(\pi^{\otimes k})
\in \mathfrak g^{(1),2}.
\end{equation}
Then:
\begin{enumerate}[label=\textup{(\alph*)}]
\item $d_\pi^2=0$;
\item $d_\pi\omega_1(\pi)=0$;
\item there is therefore a canonical cohomology class
\[
\Obs_1(\pi):=[\omega_1(\pi)]
\in H^2\bigl(\mathfrak g^{(1)},d_\pi\bigr);
\]
\item a genus-$1$ correction $\Theta^{(1)}\in\mathfrak g^{(1),1}$ exists modulo gauge if and
only if $\Obs_1(\pi)=0$;
\item if such genus-$1$ corrections exist, their gauge-equivalence classes form a torsor for
$H^1(\mathfrak g^{(1)},d_\pi)$.
\end{enumerate}
\end{theorem}

\begin{proof}
The linearization of the Maurer--Cartan equation at a Maurer--Cartan point always squares to
zero; this is the standard $L_\infty$ identity for the twisted unary bracket, proving $d_\pi^2=0$.
Now insert the ansatz
\[
\Pi(\h)=\pi+\h\Theta^{(1)}+O(\h^2)
\]
into \eqref{eq:full-modular-mc}.  The coefficient of $\h^0$ is the genus-$0$ Maurer--Cartan
equation for $\pi$, while the coefficient of $\h^1$ is exactly
\[
d_\pi\Theta^{(1)}+\omega_1(\pi)=0.
\]
Applying $d_\pi$ and using $d_\pi^2=0$ gives $d_\pi\omega_1(\pi)=0$, so
$\omega_1(\pi)$ defines a degree-$2$ cohomology class.  The displayed equation also shows that
a genus-$1$ lift exists precisely when this cocycle is exact.  Finally, if
$\Theta^{(1)}$ and $\widetilde\Theta^{(1)}$ are two solutions then their difference is
$d_\pi$-closed, and a genus-$1$ gauge transformation changes the difference by a
$d_\pi$-exact term.  This yields the torsor statement.
\end{proof}

\begin{definition}[Higher source terms]
\label{def:higher-source-terms}
Let
\[
\Pi_{<G}:=\pi+\sum_{1\le h<G}\h^h\Theta^{(h)}
\]
be a partial lift satisfying the Maurer--Cartan equation modulo $\h^G$.  The
\emph{$G$-th source term} is the coefficient
\begin{equation}
\label{eq:higher-source-term}
\omega_G(\Pi_{<G})
:=
[\h^G]\,
\sum_{r\ge 0}\sum_{k\ge 1}\frac1{k!}\,
\h^r\ell_k^{(r)}\!\bigl(\Pi_{<G}^{\otimes k}\bigr)
\in \mathfrak g^{(G),2}.
\end{equation}
\end{definition}

\begin{theorem}[All-genus obstruction tower; \ClaimStatusProvedHere]
\label{thm:all-genus-obstruction-tower}
For every $G\ge 1$, the source term $\omega_G(\Pi_{<G})$ is $d_\pi$-closed.  Hence there is a
canonical obstruction class
\[
\Obs_G(\Pi_{<G})
:=
[\omega_G(\Pi_{<G})]
\in H^2\bigl(\mathfrak g^{(G)},d_\pi\bigr).
\]
In detail:
\begin{enumerate}[label=\textup{(\roman*)}]
\item the partial lift $\Pi_{<G}$ extends to order $G$ if and only if
$\Obs_G(\Pi_{<G})=0$;
\item if extensions exist, their gauge-equivalence classes form a torsor for
$H^1(\mathfrak g^{(G)},d_\pi)$;
\item if $H^2(\mathfrak g^{(G)},d_\pi)=0$ for all $G\ge 1$, then $\pi$ admits a full modular
lift;
\item if, in addition, $H^1(\mathfrak g^{(G)},d_\pi)=0$ for all $G\ge 1$, then the modular lift
is unique up to gauge.
\end{enumerate}
\end{theorem}

\begin{proof}
Let
\[
E(\Pi):=\sum_{r\ge 0}\sum_{k\ge 1}\frac1{k!}\,\h^r\ell_k^{(r)}(\Pi^{\otimes k})
\]
be the Maurer--Cartan functional.  Since $\Pi_{<G}$ solves the MC equation modulo $\h^G$, one
has
\[
E(\Pi_{<G})=\h^G\omega_G(\Pi_{<G})+O(\h^{G+1}).
\]
For any $x\in\mathfrak g^{(G),1}$,
\[
E(\Pi_{<G}+\h^Gx)=\h^G\bigl(\omega_G(\Pi_{<G})+d_\pi x\bigr)+O(\h^{G+1}),
\]
because the only terms linear in $x$ of total order $G$ arise from genus-$0$ brackets with all
other entries equal to $\pi$.  Thus extension to order $G$ is equivalent to solving
\[
d_\pi x+\omega_G(\Pi_{<G})=0,
\]
which proves part \textup{(i)}.

To prove closure, apply the linearized Maurer--Cartan operator at $\Pi_{<G}$ to
$E(\Pi_{<G})$.  The $L_\infty$ identities imply the Bianchi-type identity
\[
d_{\Pi_{<G}}E(\Pi_{<G})=0.
\]
Reducing modulo $\h^{G+1}$ and using $d_{\Pi_{<G}}\equiv d_\pi \pmod{\h}$ yields
$d_\pi\omega_G(\Pi_{<G})=0$.

If $x$ and $x'$ are two order-$G$ lifts, then $x-x'$ is $d_\pi$-closed; order-$G$ gauge
transformations shift the difference by a $d_\pi$-exact term.  This gives part \textup{(ii)}.
The remaining statements follow by induction on $G$.
\end{proof}

\begin{corollary}[Existence and uniqueness by vanishing; \ClaimStatusProvedHere]
\label{cor:existence-uniqueness-vanishing}
If
\[
H^2\bigl(\mathfrak g^{(G)},d_\pi\bigr)=0\qquad \forall G\ge 1,
\]
then every genus-$0$ Maurer--Cartan class lifts to all genera.  If also
\[
H^1\bigl(\mathfrak g^{(G)},d_\pi\bigr)=0\qquad \forall G\ge 1,
\]
then the lift is unique up to gauge.
\end{corollary}

\subsection{Bulk--boundary comparison as a cone problem}
\label{subsec:bulk-boundary-cone-problem}

Bulk and boundary lifting problems are compared via the cone of the linearized
bulk--boundary map.

\begin{definition}[Linearized bulk--boundary map]
\label{def:linearized-bulk-boundary-map}
Let
\[
F:\mathfrak g^{\mathrm{mod}}_{\mathrm{bulk}}\longrightarrow \mathfrak g^{\mathrm{mod}}_{\partial}
\]
be a filtered $L_\infty$-morphism, and let
$\pi_{\mathrm{bulk}}\in\mathfrak g_{\mathrm{bulk}}^{(0),1}$,
$\pi_\partial\in\mathfrak g_\partial^{(0),1}$
be genus-$0$ Maurer--Cartan elements such that $F_*(\pi_{\mathrm{bulk}})$ is gauge-equivalent to
$\pi_\partial$.  The \emph{linearized bulk--boundary map} is the induced cochain map
\[
F_\pi:(\mathfrak g^{\mathrm{mod}}_{\mathrm{bulk}},d_{\pi_{\mathrm{bulk}}})
\longrightarrow
(\mathfrak g^{\mathrm{mod}}_{\partial},d_{\pi_\partial}).
\]
\end{definition}

\begin{theorem}[Relative bulk--boundary reconstruction; \ClaimStatusProvedHere]
\label{thm:relative-bulk-boundary-reconstruction}
Assume the hypotheses of Definition \ref{def:linearized-bulk-boundary-map}.  If
\[
H^1\bigl(\Cone(F_\pi)\bigr)=H^2\bigl(\Cone(F_\pi)\bigr)=0,
\]
then the genus-by-genus modular lift problem for $\pi_{\mathrm{bulk}}$ is equivalent, up to gauge,
to the genus-by-genus modular lift problem for $\pi_\partial$.  In particular:
\begin{enumerate}[label=\textup{(\roman*)}]
\item every boundary modular lift uniquely determines a compatible bulk modular lift;
\item every bulk modular lift uniquely determines a compatible boundary modular lift;
\item all extension obstructions and ambiguities are measured by the same cone cohomology.
\end{enumerate}
\end{theorem}

\begin{proof}[Proof sketch]
Compatible pairs of bulk and boundary lifts form the Maurer--Cartan space of the homotopy
fiber of $F$ over $\pi_\partial$.  The tangent complex of that homotopy fiber at the point
$(\pi_{\mathrm{bulk}},\pi_\partial)$ is $\Cone(F_\pi)[-1]$.  Standard obstruction theory for formal
moduli problems therefore places the extension obstruction in
$H^2(\Cone(F_\pi))$ and the ambiguity in $H^1(\Cone(F_\pi))$.  When both groups vanish,
compatible lifts exist and are unique at each order in $\h$, hence by induction through the
entire genus filtration.
\end{proof}

\begin{corollary}[Boundary determination of the modular bulk; \ClaimStatusProvedHere]
\label{cor:boundary-determines-bulk}
If the linearized bulk--boundary map is a quasi-isomorphism in degrees $1$ and $2$, then the
full modular bulk deformation problem is controlled by the boundary algebra.  Equivalently,
under this hypothesis the holographic dictionary is reconstructible from the boundary modular
lift alone.
\end{corollary}

\subsection{The boundary modular class}
\label{subsec:boundary-modular-class}

The genus-$1$ source term admits a canonical graph-theoretic incarnation.  This gives the
first genuinely modular invariant of the boundary algebra.

\begin{definition}[Loop-generated modular model]
\label{def:loop-generated-model}
A modular deformation package is \emph{loop-generated} if its first positive-genus operations
are generated entirely by connected stable graphs of first Betti number $1$ whose vertices are
decorated by genus-$0$ operations.  Equivalently, there are no primitive genus-$1$ vertices:
all genus-$1$ information is produced by a single loop with tree-level vertex decorations.
\end{definition}

\begin{definition}[One-wheel sum and boundary modular class]
\label{def:one-wheel-sum-and-class}
Let $(A,m_\bullet)$ be a cyclic minimal $A_\infty$-algebra, and suppose its modular
deformation package is loop-generated.  Let $\mathsf{Wheel}_1$ denote the set of connected
stable graphs of first Betti number $1$, modulo isomorphism.  For
$\Gamma\in\mathsf{Wheel}_1$, let $B_\Gamma(m_\bullet)$ be the graph contraction obtained by
placing the transferred vertices $m_n$ at the vertices of $\Gamma$ and contracting internal edges
using the cyclic pairing.  Define the completed one-wheel sum
\begin{equation}
\label{eq:one-wheel-sum}
W_1(A):=
\sum_{\Gamma\in\mathsf{Wheel}_1}^{\mathrm{comp}}
\frac{w_\Gamma}{|\Aut(\Gamma)|}
B_\Gamma(m_\bullet)
\in \mathfrak g_A^{(1),2},
\end{equation}
where $w_\Gamma$ denotes the regularized weight of the wheel type.  The cohomology class
\begin{equation}
\label{eq:boundary-modular-class}
\mathfrak W(A):=[W_1(A)]\in H^2(\mathfrak g_A^{(1)},d_{\pi_A})
\end{equation}
is called the \emph{boundary modular class} of $A$.
\end{definition}

\begin{theorem}[One-wheel reduction; \ClaimStatusProvedHere]
\label{thm:one-wheel-reduction}
Assume the modular deformation package of $(A,m_\bullet)$ is loop-generated.  Then:
\begin{enumerate}[label=\textup{(\roman*)}]
\item $W_1(A)$ is $d_{\pi_A}$-closed;
\item its cohomology class agrees with the first obstruction class,
\[
\mathfrak W(A)=\Obs_1(A,\pi_A);
\]
\item $\mathfrak W(A)$ depends only on the cyclic $A_\infty$ quasi-isomorphism class of $A$.
\end{enumerate}
\end{theorem}

\begin{proof}[Proof sketch]
In a loop-generated model, the coefficient of $\h$ in the Maurer--Cartan functional is a sum
over connected stable graphs of first Betti number $1$.  Each such graph has a unique loop.
Collapse every maximal tree attached to this loop by cyclic homological perturbation.  Each
collapsed tree becomes a single effective vertex labeled by one of the transferred operations
$m_n$.  The resulting expression is exactly the completed one-wheel sum \eqref{eq:one-wheel-sum}.
Therefore the order-$\h$ Maurer--Cartan equation reads
\[
d_{\pi_A}\Theta_A^{(1)}+W_1(A)=0,
\]
so $W_1(A)$ is closed and its cohomology class is the obstruction class.

Finally, cyclic $A_\infty$ quasi-isomorphisms induce quasi-isomorphisms of the corresponding
cyclic deformation complexes; the obstruction class of the Maurer--Cartan lifting problem is
therefore invariant under cyclic $A_\infty$ quasi-isomorphism.
\end{proof}

\begin{corollary}[Vanishing criterion and additivity; \ClaimStatusProvedHere]
\label{cor:vanishing-criterion-and-additivity}
\begin{enumerate}[label=\textup{(\roman*)}]
\item A genus-$1$ modular lift exists if and only if $\mathfrak W(A)=0$.
\item If also $H^1(\mathfrak g_A^{(1)},d_{\pi_A})=0$, then the genus-$1$ lift is unique up to
gauge.
\item If $A=A_1\oplus A_2$ is an orthogonal direct sum of cyclic minimal $A_\infty$-algebras,
then
\[
\mathfrak W(A)=\mathfrak W(A_1)\oplus\mathfrak W(A_2).
\]
\end{enumerate}
\end{corollary}

\begin{proof}
Parts \textup{(i)} and \textup{(ii)} are immediate from Theorems
\ref{thm:first-modular-obstruction} and \ref{thm:one-wheel-reduction}.  For
\textup{(iii)}, the cyclic pairing and all transferred products split block-diagonally, so every
one-wheel contraction is supported on exactly one summand.
\end{proof}

\subsection{The holomorphic BF model on $S^3$}
\label{subsec:holomorphic-bf-model-on-s3}

Specialization to Zeng's holomorphic BF model on $S^3$ \cite{ZengTwistedHolography}.

\begin{definition}[The CR-cohomological boundary algebra]
\label{def:cr-cohomological-boundary-algebra}
Let
\[
A:=\Hb = H_b^{0,0}(S^3)\oplus H_b^{0,1}(S^3).
\]
By the harmonic decomposition of the tangential Cauchy--Riemann complex and the computation
of CR cohomology, one has
\[
H_b^{0,0}(S^3)\cong \C[w_1,w_2],
\qquad
H_b^{0,1}(S^3)\cong \C[\ov w_1,\ov w_2]\,\epsilon,
\]
where $\epsilon$ is the $SU(2)$-invariant generator of $\Omega_b^{0,1}(S^3)$
\cite{ZengTwistedHolography}.  We write
\[
A^0:=\C[w_1,w_2],
\qquad
A^1:=\C[\ov w_1,\ov w_2]\,\epsilon.
\]
\end{definition}

\begin{definition}[Monomial bases and trace pairing]
\label{def:monomial-bases-and-trace-pairing}
For $p,q,r,s\ge 0$, define
\[
a_{p,q}:=w_1^p w_2^q\in A^0,
\qquad
b_{r,s}:=\frac{(r+s+1)!}{r!\,s!}\,\ov w_1^r\ov w_2^s\,\epsilon\in A^1.
\]
The cyclic trace pairing is normalized by
\[
\Tr(a_{p,q}b_{r,s})=\delta_{pr}\delta_{qs}.
\]
\end{definition}

\begin{remark}[The cyclic SDR]
\label{rem:cyclic-sdr}
Zeng constructs a special deformation retract
\[
h\,(\Omega_b^{0,\bullet}(S^3),\delbar_b)
\underset{i}{\overset{p}{\rightleftarrows}}
(\HCR,0)
\]
with $h\circ i=0$, $p\circ h=0$, $h^2=0$, and $\id-ip=\delbar_b h+h\delbar_b$.
The retract is compatible with the cyclic structure, and therefore transfers the
commutative dg algebra structure on $\Omega_b^{0,\bullet}(S^3)$ to a cyclic $A_\infty$ (indeed
$C_\infty$) structure $m_\bullet=\{m_n\}_{n\ge 2}$ on $A$ \cite{ZengTwistedHolography}.
This is the transferred boundary algebra governing the KK interactions.
\end{remark}

\subsection{A concrete one-vertex calculus}
\label{subsec:concrete-one-vertex-calculus}

Cyclic traces of the transferred operations carry the first modular data.  The one-vertex sectors have a rigid arithmetic structure.

\begin{proposition}[The $m_2$ shift formula; \ClaimStatusProvedHere]
\label{prop:m2-shift-formula}
For all $p,q,r,s\ge 0$,
\[
m_2(a_{p,q},b_{r,s})=
\begin{cases}
 b_{r-p,s-q},& p\le r,\ q\le s,\\[4pt]
 0,& \text{otherwise.}
\end{cases}
\]
\end{proposition}

\begin{proof}
Zeng's explicit formula for the transferred binary product is
\[
m_2\!\left(w_1^p w_2^q,
\frac{(r+s+1)!}{r!s!}\,\ov w_1^r \ov w_2^s\epsilon\right)
=
\frac{(r+s-p-q+1)!}{(r-p)!(s-q)!}\,\ov w_1^{r-p}\ov w_2^{s-q}\epsilon.
\]
Recognizing the right-hand side as $b_{r-p,s-q}$ proves the claim.
\end{proof}

\begin{definition}[One-vertex traced sectors]
\label{def:one-vertex-traced-sectors}
For $n\ge 2$ and $\beta_1,\dots,\beta_{n-1}\in A^1$, define the traced one-vertex sector
\begin{equation}
\label{eq:one-vertex-traced-sector}
Q_n(\beta_1,\dots,\beta_{n-1})
:=
\sum_{p,q\ge 0}
\Tr\!\bigl(a_{p,q}\,m_n(a_{p,q},\beta_1,\dots,\beta_{n-1})\bigr).
\end{equation}
For $n=2$, write
\[
\theta_{r,s}:=Q_2(b_{r,s}).
\]
\end{definition}

\begin{theorem}[The tadpole projector; \ClaimStatusProvedHere]
\label{thm:tadpole-projector}
The coefficients $\theta_{r,s}$ satisfy
\[
\theta_{r,s}=
\begin{cases}
1,& r,s\ \text{both even},\\
0,& \text{otherwise}.
\end{cases}
\]
Equivalently,
\[
W^{(2)}_{\mathrm{alg}}:=\sum_{r,s\ge 0}\theta_{r,s}b_{r,s}
=\sum_{m,n\ge 0} b_{2m,2n}.
\]
In particular, the quadratic one-vertex wheel is the projector onto the even--even KK lattice.
\end{theorem}

\begin{proof}
By Proposition \ref{prop:m2-shift-formula},
\[
\theta_{r,s}
=
\sum_{p,q\ge 0}\Tr\bigl(a_{p,q} b_{r-p,s-q}\bigr)
=
\sum_{p,q\ge 0}\delta_{p,r-p}\delta_{q,s-q}.
\]
This has a solution if and only if $2p=r$ and $2q=s$, i.e. if and only if $r$ and $s$ are both
even.  When a solution exists it is unique, so the sum equals $1$.
\end{proof}

\begin{corollary}[Lowest-mode tadpole; \ClaimStatusProvedHere]
\label{cor:lowest-mode-tadpole}
Let
\[
L:=\Span\{b_{0,0},b_{1,0},b_{0,1}\}\subset A^1.
\]
Then the tadpole projector restricts to
\[
W^{(2)}_{\mathrm{alg}}\big|_L=b_{0,0}.
\]
If one restores the geometric genus-$1$ weight $w_2$, the physical tadpole truncation on the
lowest modes is $w_2\,b_{0,0}$.
\end{corollary}

\begin{theorem}[One-vertex trace selection rule; \ClaimStatusProvedHere]
\label{thm:one-vertex-trace-selection-rule}
Let $n\ge 2$, and let $\beta_i=b_{r_i,s_i}\in A^1$ for $1\le i\le n-1$.  Set
\[
R:=\sum_{i=1}^{n-1} r_i,
\qquad
S:=\sum_{i=1}^{n-1} s_i.
\]
Then for every $p,q\ge 0$ one has
\[
m_n(a_{p,q},b_{r_1,s_1},\dots,b_{r_{n-1},s_{n-1}})
\in
\C\cdot b_{R-p+n-2,\,S-q+n-2}.
\]
Consequently,
\[
Q_n(b_{r_1,s_1},\dots,b_{r_{n-1},s_{n-1}})\neq 0
\]
can occur only if
\begin{equation}
\label{eq:selection-rule}
2p=R+n-2,
\qquad
2q=S+n-2.
\end{equation}
In particular, for fixed external labels there is at most one internal KK mode contributing to
the one-vertex trace.
\end{theorem}

\begin{proof}
Zeng's explicit formula for $m_n$ in the orthonormal $SU(2)$ basis shows that the output of
\[
m_n\bigl(e_{m_0}^{(j_0)},\ov e_{m_1}^{(\ov j_1)}\epsilon,\dots,
\ov e_{m_{n-1}}^{(\ov j_{n-1})}\epsilon\bigr)
\]
lies in the unique representation
\[
\ov e_{M_{n-1}}^{(\ov J_{n-1}-j_0+n-2)}\epsilon,
\qquad
\ov J_{n-1}=\sum_{i=1}^{n-1}\ov j_i,
\qquad
M_{n-1}=m_0+\sum_{i=1}^{n-1}m_i.
\]
Translate this to monomial labels by
\[
j_0=\frac{p+q}{2},\qquad m_0=\frac{p-q}{2},
\qquad
\ov j_i=\frac{r_i+s_i}{2},\qquad m_i=\frac{s_i-r_i}{2}.
\]
Writing the output as $b_{u,v}$ means
\[
\frac{u+v}{2}=\ov J_{n-1}-j_0+n-2,
\qquad
\frac{v-u}{2}=M_{n-1}.
\]
Solving for $(u,v)$ gives
\[
u=R-p+n-2,
\qquad
v=S-q+n-2.
\]
Tracing against $a_{p,q}$ forces $(u,v)=(p,q)$, which is exactly
\eqref{eq:selection-rule}.  Therefore there is at most one internal mode.
\end{proof}

\begin{theorem}[Low-mode parity law; \ClaimStatusProvedHere]
\label{thm:low-mode-parity-law}
Restrict the one-vertex sectors to the lowest-mode space
\[
L=\Span\{b_{0,0},b_{1,0},b_{0,1}\}.
\]
For a tensor in $L^{\otimes(n-1)}$, let $N_{10}$ denote the number of $b_{1,0}$ inputs and
$N_{01}$ the number of $b_{0,1}$ inputs.  Then
\[
Q_n\big|_{L^{\otimes(n-1)}}\neq 0
\qquad\Longrightarrow\qquad
N_{10}\equiv N_{01}\equiv n\pmod 2.
\]
\end{theorem}

\begin{proof}
On $L$, one has $R=N_{10}$ and $S=N_{01}$.  By Theorem
\ref{thm:one-vertex-trace-selection-rule}, nonvanishing requires
$R+n-2$ and $S+n-2$ to be even.  Since $n-2\equiv n\pmod 2$, the statement follows.
\end{proof}

\begin{corollary}[Cubic support on the lowest modes; \ClaimStatusProvedHere]
\label{cor:cubic-support-lowest-modes}
For $n=3$, the only nonzero low-mode one-vertex channels are the mixed ordered pairs
\[
(b_{1,0},b_{0,1})
\qquad\text{and}\qquad
(b_{0,1},b_{1,0}).
\]
No ordered pair involving $b_{0,0}$ contributes to $Q_3|_{L^{\otimes 2}}$.
\end{corollary}

\begin{proof}
For $n=3$, Theorem \ref{thm:low-mode-parity-law} requires
$N_{10}\equiv N_{01}\equiv 1\pmod 2$.  Since there are only two inputs, the only possibility is
one copy of $b_{1,0}$ and one copy of $b_{0,1}$.
\end{proof}

\begin{proposition}[Vanishing of the symmetric cubic one-vertex sector; \ClaimStatusProvedHere]
\label{prop:vanishing-symmetric-cubic-one-vertex-sector}
The symmetric cubic one-vertex sector vanishes on the lowest modes.  Equivalently, the
trilinear form $Q_3|_{L^{\otimes 2}}$ factors through an alternating $2$-form
\[
\omega_{\mathrm{pm}}\in \Lambda^2 L^*,
\]
which we call the \emph{low-mode pre-modular two-form}.
\end{proposition}

\begin{proof}
By Corollary \ref{cor:cubic-support-lowest-modes}, the only nonzero channels are the two mixed
ordered pairs.  Zeng's antisymmetry relation for the cubic coefficients shows that swapping the
two $A^1$ inputs reverses the coefficient.  Therefore the symmetric combination cancels, and the
result factors through $\Lambda^2L^*$.
\end{proof}

\begin{theorem}[Quartic support theorem on the lowest modes; \ClaimStatusProvedHere]
\label{thm:quartic-support-theorem-lowest-modes}
For $n=4$, the one-vertex low-mode sector is supported on exactly three unordered channels:
\[
(b_{0,0},b_{0,0},b_{0,0}),
\qquad
(b_{1,0},b_{1,0},b_{0,0}),
\qquad
(b_{0,1},b_{0,1},b_{0,0}).
\]
Equivalently, on $\Sym^3(L)$ one has a unique decomposition
\begin{equation}
\label{eq:q4-normal-form}
Q_4\big|_{\Sym^3(L)}
=
Q_{000}\,b_{0,0}^3
+
Q_{110}\,b_{1,0}^2b_{0,0}
+
Q_{011}\,b_{0,1}^2b_{0,0},
\end{equation}
for uniquely defined scalars $Q_{000},Q_{110},Q_{011}\in\C$.
\end{theorem}

\begin{proof}
For $n=4$, Theorem \ref{thm:low-mode-parity-law} gives
$N_{10}\equiv N_{01}\equiv 0\pmod 2$.  Since there are three inputs,
\[
N_{10}+N_{01}+N_{00}=3,
\]
so the only possible triples are
\[
(N_{10},N_{01},N_{00})=(0,0,3),\ (2,0,1),\ (0,2,1).
\]
These are exactly the three unordered channels above.
\end{proof}

\begin{corollary}[The charged quartic sector is two-dimensional; \ClaimStatusProvedHere]
\label{cor:charged-quartic-sector-two-dimensional}
After removing the neutral channel $b_{0,0}^3$, the first charged one-vertex quartic sector on
low modes lies in the $2$-dimensional space
\[
\Span\{b_{1,0}^2b_{0,0},\ b_{0,1}^2b_{0,0}\}.
\]
Thus the first charged quartic modular datum is encoded by the pair $(Q_{110},Q_{011})$.
\end{corollary}

\begin{remark}[Tiny internal support]
\label{rem:tiny-internal-support}
The trace selection rule determines the internal test mode uniquely.  For the three quartic
channels of Theorem \ref{thm:quartic-support-theorem-lowest-modes}, one finds
\[
Q_{000}: \ a_{1,1},
\qquad
Q_{110}: \ a_{2,1},
\qquad
Q_{011}: \ a_{1,2}.
\]
The charged quartic coefficients are therefore algebraically finite, without regularization.  By contrast, generic loop sectors of the KK theory require infinite KK sums \cite{ZengTwistedHolography}.
\end{remark}

\subsection{The first two-vertex bubble and the exceptional finite locus}
\label{subsec:first-two-vertex-bubble-exceptional-finite-locus}

The first loop-shaped graph is the two-vertex cubic--cubic wheel.  The lowest-mode sector remains rigid.

\begin{theorem}[Support of the first mixed cubic--cubic bubble; \ClaimStatusProvedHere]
\label{thm:support-first-mixed-cubic-cubic-bubble}
Consider the first two-vertex wheel built from two cubic vertices.  Restricted to the lowest-mode
subspace $L$, its symmetric external content is supported on a single monomial:
\[
b_{1,0}^2 b_{0,1}^2.
\]
No monomial involving $b_{0,0}$ occurs in this first cubic--cubic bubble sector.
\end{theorem}

\begin{proof}
By Corollary \ref{cor:cubic-support-lowest-modes}, each cubic vertex is nonzero only if its two
$A^1$ inputs are one copy of $b_{1,0}$ and one copy of $b_{0,1}$.  Therefore every nonzero
cubic--cubic wheel contributes exactly one $b_{1,0}$ and one $b_{0,1}$ from each vertex.  The
total external monomial is thus forced to be $b_{1,0}^2b_{0,1}^2$.
\end{proof}

\begin{theorem}[Exceptional finiteness of the first mixed bubble; \ClaimStatusProvedHere]
\label{thm:exceptional-finiteness-first-mixed-bubble}
The first mixed cubic--cubic bubble on the lowest modes is finite before regularization.  More
precisely, for the mixed cubic sector one has
\[
m_3(a_{p,q},b_{1,0},b_{0,1})\in \C\cdot b_{2-p,2-q},
\]
and therefore only the finite range
\[
0\le p,q\le 2
\]
can contribute to the first mixed bubble coefficient.  In particular, the first mixed bubble is a
finite algebraic contraction of finitely many cubic structure constants.
\end{theorem}

\begin{proof}
Apply Theorem \ref{thm:one-vertex-trace-selection-rule} with $n=3$ and external labels
$(r_1,s_1)=(1,0)$, $(r_2,s_2)=(0,1)$.  Then $R=S=1$, so the cubic output is
\[
b_{R-p+1,S-q+1}=b_{2-p,2-q}.
\]
For this to make sense in $A^1$, one must have $2-p\ge 0$ and $2-q\ge 0$, hence
$0\le p,q\le 2$.  Thus the internal mode sum is finite.
\end{proof}

\begin{definition}[The low-mode charged modular jet]
\label{def:low-mode-charged-modular-jet}
The \emph{low-mode charged modular jet} of the holomorphic BF boundary algebra is the triple
\begin{equation}
\label{eq:charged-modular-jet}
\mathfrak J_{\mathrm{ch}}^{\mathrm{low}}:=(Q_{110},Q_{011},B_{\mathrm{mix}}),
\end{equation}
where $Q_{110}$ and $Q_{011}$ are the charged one-vertex quartic coefficients of
\eqref{eq:q4-normal-form}, and $B_{\mathrm{mix}}$ is the coefficient of
$b_{1,0}^2b_{0,1}^2$ in the first cubic--cubic bubble sector.
\end{definition}

\begin{theorem}[Low-mode modular normal form; \ClaimStatusProvedHere]
\label{thm:low-mode-modular-normal-form}
Up to and including the first cubic--cubic bubble order, the symmetric modular shadow of the
holomorphic BF boundary algebra on the lowest-mode sector has the form
\begin{equation}
\label{eq:low-mode-modular-normal-form}
W_{\mathrm{low}}
=
w_2\,b_{0,0}
+
Q_{000}\,b_{0,0}^3
+
Q_{110}\,b_{1,0}^2b_{0,0}
+
Q_{011}\,b_{0,1}^2b_{0,0}
+
B_{\mathrm{mix}}\,b_{1,0}^2b_{0,1}^2
+
\textup{(higher sectors)}.
\end{equation}
No other symmetric low-mode monomial can appear at this order.
\end{theorem}

\begin{proof}
The first term is the tadpole contribution from Corollary
\ref{cor:lowest-mode-tadpole}.  The one-vertex cubic symmetric contribution vanishes by
Proposition \ref{prop:vanishing-symmetric-cubic-one-vertex-sector}.  The one-vertex quartic
contribution is exhausted by Theorem \ref{thm:quartic-support-theorem-lowest-modes}.  The
first two-vertex bubble contribution is exhausted by Theorem
\ref{thm:support-first-mixed-cubic-cubic-bubble}.  No further symmetric low-mode monomial can
occur before higher-valence or higher-vertex sectors enter.
\end{proof}

\begin{remark}[What has been proved]
\label{rem:what-has-been-proved}
Theorem \ref{thm:low-mode-modular-normal-form} is the sharpest rigorous reduction obtained
in this section.  It does \emph{not} compute the three charged numbers
$Q_{110},Q_{011},B_{\mathrm{mix}}$, but it proves that these are the \emph{only} charged low-mode
symmetric modular coefficients that can appear at the first nontrivial orders.  This converts an
infinite KK/modular problem into a finite jet problem.
\end{remark}

\subsection{A strengthened frontier: new structures and new targets}
\label{subsec:strengthened-frontier-new-structures}

The low-mode reduction yields the following invariants.

\begin{definition}[The low-mode pre-modular package]
\label{def:low-mode-pre-modular-package}
The \emph{low-mode pre-modular package} of the holomorphic BF boundary algebra is the pair
\[
\bigl(\omega_{\mathrm{pm}},\mathfrak J_{\mathrm{ch}}^{\mathrm{low}}\bigr),
\]
where $\omega_{\mathrm{pm}}\in \Lambda^2L^*$ is the alternating cubic one-vertex shadow of
Proposition \ref{prop:vanishing-symmetric-cubic-one-vertex-sector}, and
$\mathfrak J_{\mathrm{ch}}^{\mathrm{low}}$ is the charged modular jet of Definition
\ref{def:low-mode-charged-modular-jet}.
\end{definition}

\begin{proposition}[Functoriality under cyclic equivalence on the lowest modes; \ClaimStatusProvedHere]
\label{prop:functoriality-lowest-modes}
Let $A$ and $A'$ be cyclic minimal $A_\infty$ models whose lowest-mode sectors are identified
by a cyclic $A_\infty$ quasi-isomorphism preserving the subspace
$L=\Span\{b_{0,0},b_{1,0},b_{0,1}\}$.  Then:
\begin{enumerate}[label=\textup{(\roman*)}]
\item the low-mode pre-modular two-form $\omega_{\mathrm{pm}}$ is carried to the corresponding
$\omega_{\mathrm{pm}}'$;
\item the support constraints of Theorems
\ref{thm:quartic-support-theorem-lowest-modes} and
\ref{thm:support-first-mixed-cubic-cubic-bubble} are preserved;
\item the charged modular jet transforms functorially.
\end{enumerate}
\end{proposition}

\begin{proof}
The support theorems depend only on the transferred $A_\infty$ structure, the cyclic pairing,
and the grading/weight bookkeeping encoded by the selection rule.  A cyclic
$A_\infty$ quasi-isomorphism preserving $L$ transports all of this data.  Therefore the induced
trilinear and quartic forms on $L$ and the first bubble coefficient transform functorially.
\end{proof}

\begin{problem}[The quartic Wigner identity problem]
\label{prob:quartic-wigner-identity}
Find a closed-form representation-theoretic identity, directly in terms of Clebsch--Gordan and
Wigner-symbol data, for the charged quartic coefficients
\[
Q_{110}=\Tr\bigl(a_{2,1}\,m_4(a_{2,1},b_{1,0},b_{1,0},b_{0,0})\bigr),
\qquad
Q_{011}=\Tr\bigl(a_{1,2}\,m_4(a_{1,2},b_{0,1},b_{0,1},b_{0,0})\bigr).
\]
The support theory shows that these are finite, tiny-support Wigner sums.
A closed-form evaluation remains open.
\end{problem}

\begin{problem}[Universal finite-locus theorem]
\label{prob:universal-finite-locus}
Characterize, in an arbitrary loop-generated cyclic $A_\infty$ model, exactly which low-mode
wheel sectors remain finite before regularization.  The holomorphic BF calculation exhibits an \emph{exceptional finite locus}: a representation-theoretically defined region of
external weight space on which loop diagrams collapse to finite algebraic contractions, while generic loop sectors diverge.
\end{problem}

\begin{problem}[The modular bootstrap problem]
\label{prob:modular-bootstrap}
Use the quadratic identities of the BRST anomaly/higher-operation formalism to bootstrap the
first modular boundary coefficients from associativity and stable-graph factorization,
rather than direct Feynman summation.  The target is a boundary modular analog of
the bootstrap constraints on celestial one-loop OPE data.
\end{problem}

\begin{problem}[Quantization of the boundary modular class]
\label{prob:quantization-boundary-modular-class}
Relate the genus-$1$ boundary modular class and the charged modular jet to deformation
quantization of the underlying coisson/Poisson-vertex structure.  The bridge is the
$3$d holomorphic--topological Poisson-vertex framework of Khan--Zeng and the
configuration-space/Feynman realization of higher operations and their quadratic identities.
\end{problem}


\section{Celestial OPE from the bar complex: gluon and graviton sectors}
\label{sec:celestial-ope-bar-complex}

The celestial holography programme maps $4$-dimensional scattering amplitudes to $2$-dimensional celestial CFT correlators via Mellin transform.  The boundary chiral algebras of the holomorphic-topological framework (affine Kac--Moody for gauge theory, Virasoro for gravity) encode celestial OPE data through their bar complexes.  The collision residues of the bar differential $d_{\barBch}$ extract these OPE coefficients via the $d\log$ absorption mechanism (AP19): the bar kernel $d\log(z_1 - z_2)$ absorbs one pole order, so an OPE singularity $\sim z^{-N}$ produces an $r$-matrix pole $\sim z^{-(N-1)}$.

\subsection{Gluon celestial OPE from $\barBch(\widehat{\fg}_k)$}
\label{subsec:ch-core-gluon-celestial-ope}

The holomorphic-topological twist of $4$d $\mathcal{N} = 2$ Yang--Mills with gauge algebra $\fg$ on $\mathbb{C} \times \mathbb{R}_{\ge 0}$ produces a boundary Kac--Moody algebra $\widehat{\fg}_k$.  The positive-helicity gluon maps to the affine current $J^a(z)$.  The $\lambda$-bracket $\{J^a {}_\lambda J^b\} = f^{ab}{}_c\, J^c + k\,\kappa^{ab}\,\lambda$ Laplace-transforms to $r^{\mathrm{Lap}}(z) = f^{ab}{}_c\, J^c / z + k\,\kappa^{ab}/z^2$.  After $d\log$ absorption, the $z^{-1}$ structure-constant term drops to $z^0$ and the $z^{-2}$ level term becomes $z^{-1}$:
\[
r^{\mathrm{KM}}(z) \;=\; \frac{k\,\Omega}{z}
\;=\; \frac{k\,\kappa^{ab}\,t_a \otimes t_b}{z},
\]
where $\Omega$ is the Casimir tensor.

\begin{evidence}[Celestial gluon OPE from the bar complex; \ClaimStatusHeuristic]
\label{ev:ch-core-celestial-gluon-ope}
Let $\mathcal{O}^{+,a}_{\Delta}(z)$ denote the celestial primary of a positive-helicity gluon of colour~$a$ and conformal dimension~$\Delta$.  The celestial OPE extracted from the collision residue of $\barBch(\widehat{\fg}_k)$ is
\[
\mathcal{O}^{+,a}_{\Delta_1}(z_1)\,\mathcal{O}^{+,b}_{\Delta_2}(z_2)
\;\sim\;
\frac{f^{ab}{}_c}{z_1 - z_2}\;
\mathcal{O}^{+,c}_{\Delta_1 + \Delta_2 - 1}(z_2)
\;+\; \textup{regular},
\]
with OPE coefficient the structure constant $f^{ab}{}_c$ of the gauge Lie algebra~$\fg$.
The collision $r$-matrix $r^{\mathrm{KM}}(z) = k\,\Omega/z$ encodes the braiding of two line operators at spectral distance~$z$.  In the colour-ordered projection, $\langle c|\,\Omega\,|a,b\rangle = f^{ab}{}_c$, producing the simple-pole OPE $\sim f^{ab}{}_c/(z_1 - z_2)$.  The conformal-dimension shift $\Delta_1 + \Delta_2 - 1$ arises from the Mellin convolution beta-function integral $B(\Delta_1, \Delta_2) = \Gamma(\Delta_1)\,\Gamma(\Delta_2)/\Gamma(\Delta_1 + \Delta_2)$.  The level~$k$ drops out because it is absorbed into the normalisation of celestial primaries.  The argument is heuristic: it assumes the celestial primary basis diagonalises the bar-complex collision residue; a full proof requires the Mellin-space factorisation theorem for the holomorphic-topological propagator.
\end{evidence}

\begin{remark}[Comparison with Fan--Fotopoulos--Taylor]
\label{rem:ch-core-fft}
This matches the celestial gluon OPE obtained by direct Mellin transform of the Parke--Taylor amplitude.  Here $f^{ab}{}_c$ is extracted from the collision residue of the bar complex via the $d\log$ absorption mechanism, not inserted by hand.
\end{remark}


\subsection{Graviton celestial OPE from $\barBch(\mathrm{Vir}_c)$}
\label{subsec:ch-core-graviton-celestial-ope}

The holomorphic-topological twist of $3$d gravity on $\mathbb{C} \times \mathbb{R}_{\ge 0}$ produces a boundary Virasoro algebra $\mathrm{Vir}_c$.  The positive-helicity graviton maps to the stress tensor~$T(z)$.  The Virasoro $\lambda$-bracket $\{T{}_\lambda T\} = (\partial + 2\lambda)\,T + \frac{c}{12}\,\lambda^3$ Laplace-transforms to $r^{\mathrm{Lap}}(z) = \partial T/z + 2T/z^2 + (c/2)/z^4$.  After $d\log$ absorption (AP19), the $\partial T$ term drops, producing the two-pole collision $r$-matrix:
\[
r^{\mathrm{Vir}}(z) \;=\; \frac{c/2}{z^3} \;+\; \frac{2T}{z}.
\]

\begin{evidence}[Celestial graviton OPE from the bar complex; \ClaimStatusHeuristic]
\label{ev:ch-core-celestial-graviton-ope}
The celestial OPE extracted from the collision residue of $\barBch(\mathrm{Vir}_c)$ is
\[
\mathcal{O}^+_{\Delta_1}(z_1)\,\mathcal{O}^+_{\Delta_2}(z_2)
\;\sim\;
\frac{\Delta_2 - 1}{z_1 - z_2}\;
\mathcal{O}^+_{\Delta_1 + \Delta_2 - 1}(z_2)
\;+\;
\frac{c/2}{(z_1 - z_2)^3}\;
\mathcal{O}^+_{\Delta_1 + \Delta_2 - 3}(z_2)
\;+\; \textup{regular}.
\]
The leading coefficient $(\Delta_2 - 1)$ arises from $r_1 = 2T$ (the Weinberg soft graviton theorem); the subleading coefficient $c/2$ arises from $r_3 = c/2$ (the gravitational anomaly).  The two-pole structure matches the Pasterski--Shao--Strominger graviton OPE.
\end{evidence}

\begin{proof}[Evidence]
The pole $r_1 = 2T$ acts on a celestial primary of dimension $\Delta_2$ via the Virasoro zero mode, yielding the coefficient $(\Delta_2 - 1)$ after the spin-$2$ conformal weight shift.  The pole $r_3 = c/2$ is a central element producing the triple-pole term with output dimension $\Delta_1 + \Delta_2 - 3$.  The absence of a $z^{-2}$ pole is the bar-complex manifestation of the $d\log$ absorption of the $\partial T$ term.  No higher poles exist because the Virasoro $\lambda$-bracket is cubic in~$\lambda$.  As with the gluon case, the argument is heuristic: the extraction of celestial OPE coefficients from the bar-complex collision residue assumes the Mellin-space factorisation of the holomorphic-topological propagator, which has not been proved in full generality.
\end{proof}


\subsection{Helicity-chirality correspondence}
\label{subsec:ch-core-helicity-chirality}

The bar complex sees exactly one helicity.  The two colours of $\SCchtop$ are the two helicities, and the no-open-to-closed rule is the helicity selection rule (Observation~\ref{obs:helicity-directionality}).

\begin{theorem}[Helicity splitting via the Swiss-cheese decomposition; part~\textup{(i)} \ClaimStatusProvedHere, parts~\textup{(ii)--(v)} \ClaimStatusHeuristic]
\label{thm:ch-core-helicity-splitting}
In the holomorphic-topological framework:
\begin{enumerate}[label=\textup{(\roman*)}]
\item The bar complex $\barBch(\mathcal{A})$ is built from logarithmic forms on $\overline{\mathrm{FM}}_n(\mathbb{C})$, which are holomorphic.  It encodes exactly the self-dual (positive-helicity) sector.
\item The Koszul dual bar complex $\barBch(\mathcal{A}^!)$ encodes the anti-self-dual (negative-helicity) sector: Koszul duality $\mathcal{A} \mapsto \mathcal{A}^!$ exchanges holomorphic and anti-holomorphic collision data.
\item Mixed-helicity amplitudes arise from the pairing between $\barBch(\mathcal{A})$ and $\barBch(\mathcal{A}^!)$ mediated by the universal twisting morphism $\tau \colon \barBch(\mathcal{A}) \to \mathcal{A}^!$.  Physically, this pairing is the bulk propagator $K(z,\bar z,t) \sim \Theta(t)/(2\pi z)$ connecting holomorphic and anti-holomorphic insertions.
\item All-positive and single-negative tree amplitudes vanish: the weight form $\omega_n^{(+,\ldots,+)}$ has the wrong form-type to pair nontrivially with the volume form on $\overline{\mathrm{FM}}_n(\mathbb{C})$.
\item MHV amplitudes (two negative helicities) are the simplest non-vanishing pairing between $\barBch(\mathcal{A})$ and $\barBch(\mathcal{A}^!)$: the Parke--Taylor formula encodes the structure of this minimal pairing.
\end{enumerate}
In the Swiss-cheese language: the closed colour (holomorphic, $\mathbb{C}$-direction) encodes positive helicity via the bar differential; the open colour (topological, $\mathbb{R}$-direction) encodes negative helicity via the Koszul dual.  Mixed-helicity amplitudes are open--closed couplings under $\mathrm{SC}^{\mathrm{ch,top}}$.  The no-open-to-closed rule corresponds to the statement that negative-helicity states cannot generate positive-helicity states by OPE alone.

\emph{Helicity is chirality; helicity splitting is the Swiss-cheese decomposition.}
\end{theorem}

\begin{proof}[Proof of part~\textup{(i)}]
The bar differential $d_{\barBch} = d_Q + d_{\mathrm{res}} + d_{A_\infty}$ is built from collision residues of the holomorphic OPE on $\overline{\mathrm{FM}}_n(\mathbb{C})$.  The logarithmic forms $\Omega^\bullet_{\log}(\overline{\mathrm{FM}}_n(\mathbb{C}))$ are generated by $d\log(z_i - z_j)$ and carry no $d\bar z$ factors: they span a pure Hodge type~$(p,0)$ within the mixed Hodge structure on $H^\bullet(\overline{\mathrm{FM}}_n(\mathbb{C}))$.  Every element of $\barBch(\mathcal{A})$ is therefore a purely holomorphic object: it sees exactly the self-dual sector and nothing else.
\end{proof}

\begin{evidence}[Justification of parts~\textup{(ii)--(v)}]
The following arguments support the remaining identifications but do not constitute complete proofs.  Rigorous verification of parts~(ii)--(v) requires a full analysis of the form-type bookkeeping under Verdier duality and the Koszul pairing on $\overline{\mathrm{FM}}_n(\mathbb{C})$.

\emph{(ii)} The Koszul dual $\mathcal{A}^! = H^\bullet(\mathbb{D}_{\mathrm{Ran}}\,\barBch(\mathcal{A}))$ (Verdier duality, not cobar) carries a chiral structure whose bar complex encodes the dual collision structure, the anti-holomorphic OPE.  The Koszul involution exchanges holomorphic and anti-holomorphic data because it dualises the weight-grading on the generating space~$V$.  The identification of this dual collision structure with the negative-helicity sector rests on the physical interpretation of Verdier duality as CPT conjugation, which reverses helicity.

\emph{(iii)} The twisting morphism $\tau \colon \barBch(\mathcal{A}) \to \mathcal{A}^!$ is the universal coupling classified by the bar complex.  Its Mellin transform produces the mixed-helicity vertex.  The identification with the bulk propagator $K(z,\bar z,t)$ is physically motivated but requires a rigorous Mellin-space analysis to verify the precise coefficient matching.

\emph{(iv)} Form-type obstruction: $\omega_n^{(+,\ldots,+)}$ is purely holomorphic and has vanishing integral over $\overline{\mathrm{FM}}_n(\mathbb{C})$ by type.  The one-minus amplitude vanishes because a single anti-holomorphic insertion cannot supply the required $(0,n{-}3)$-form degree needed for a non-degenerate pairing.

\emph{(v)} Two anti-holomorphic insertions provide the minimum numerator $\langle ij\rangle^4$ via the anti-holomorphic Penrose transform; the positive-helicity insertions supply $\prod\langle k,k{+}1\rangle^{-1}$ via the holomorphic $d\log$ residues.  The $n = 4$ case is verified rigorously: the bar-complex integral over $\overline{\mathrm{FM}}_4(\mathbb{C})$ reproduces the colour-ordered MHV amplitude (Theorem~\ref{thm:cbt-mhv-bar-integral}).
\end{evidence}

\begin{remark}[Genus $g \geq 1$ persistence of helicity splitting]
\label{rem:ch-core-helicity-splitting-genus}
At genus $g \geq 1$, the curved bar differential satisfies $d^2 = \kappa(\mathcal{A}) \cdot \omega_g$ where $\omega_g$ is the Arakelov-corrected Hodge class.  Although this curvature involves non-holomorphic data (the Arakelov Green's function depends on both $z$ and $\bar z$ through the modular parameter~$\tau$), it enters as a \emph{central scalar} in the algebra; it does not introduce anti-holomorphic forms into the bar complex elements themselves.  The helicity splitting therefore persists at all genera: the curvature mixes holomorphic and anti-holomorphic data only through the modular parameter~$\tau$, not through the $z$-coordinates on the celestial sphere.  In physical terms, the genus-$g$ correction deforms the \emph{coupling constants} of the self-dual sector (via the curved $A_\infty$ structure) but does not mix the self-dual and anti-self-dual sectors.
\end{remark}

\begin{remark}[Colour-ordered amplitudes and the two-colour architecture]
\label{rem:ch-core-colour-ordered-two-colour}
The colour-ordered celestial amplitudes appearing in
Theorem~\ref{thm:ch-core-helicity-splitting}(v) are
intrinsically \emph{ordered-bar} objects: they live on ordered
configurations $\Conf_n^{<}(\C)$, not on the symmetric quotient
$\Conf_n(\C)/\Sigma_n$.  Full (colour-summed) amplitudes are
recovered by $R$-matrix descent
(Proposition~\ref{prop:r-matrix-descent}):
\[
  \mathcal{M}_n
  \;=\;
  \sum_{\sigma \in \Sigma_n / \mathrm{cyc}}
  R_\sigma \cdot \mathcal{M}_n^{\mathrm{ord}}(\sigma),
\]
where $R_\sigma$ is the monodromy of the $R$-matrix along the
path permuting the insertion points to the ordering~$\sigma$.
This is the celestial manifestation of the two-colour
architecture developed in Part~\ref{part:e1-core}: the ordered
bar complex $\Barchord(\cA)$ carries the colour-ordered
data, the $R$-matrix provides the descent datum, and the
unordered bar complex $\Barch(\cA)$ carries the full amplitudes
as $R$-twisted $\Sigma_n$-coinvariants
(Chapter~\ref{ch:ordered-associative-chiral-kd},
\S\ref{sec:ordered-bar-explicit}).  The Parke--Taylor denominator
$\prod_{i=1}^n \langle i,\, i{+}1 \rangle^{-1}$ is the
$d\log$ residue on $\Conf_n^{<}(\C)$, the holomorphic colour
of the ordered bar differential, while the colour-sum
over~$\Sigma_n/\mathrm{cyc}$ is the topological colour's descent.
\end{remark}

\begin{remark}[Conformally soft theorems as bar cohomology]
\label{rem:ch-core-soft-bar-cohomology}
\index{soft graviton theorem!bar cohomology}
\index{celestial holography!soft theorems}
The soft graviton theorem at sub${}^n$-leading order is a bar
cohomology class of the $\mathfrak{w}_{1+\infty}$ bar complex
at conformal dimension $\Delta = 1 - n$.  The hierarchy:
$S_0$ (supertranslation / BMS), $S_1$ (superrotation /
Virasoro extension of Lorentz), $S_2$ ($\mathfrak{w}_{1+\infty}$
extension / spin-$3$ soft theorem).  For self-dual Yang--Mills
with gauge group~$G$, the collinear algebra at level $k = 0$
is chirally Koszul: $H^*(\Barch(J_{\mathrm{col}}))$ is concentrated
in bar degree~$1$, giving the Koszul duality
$\mathrm{Com}^! = \mathrm{Lie}$.
The collision residue $r_{ab}(z) = f^c_{ab}/z$ is the celestial
R-matrix (AP19: simple pole from the simple-pole OPE via
$d\log$ extraction).
Computational verification:
\texttt{compute/lib/celestial\_koszul\_ope.py}.
\end{remark}

\begin{remark}[$\mathfrak{w}_{1+\infty}$ shadow obstruction tower and celestial genus expansion]
\label{rem:ch-core-w-shadow}
\index{celestial holography!shadow tower}
The modular characteristic and higher shadow coefficients of
$\mathfrak{w}_{1+\infty}[\lambda = 1]$ control the genus
expansion of celestial amplitudes.  On the T-line (Virasoro
sub-tower), the shadow data is that of $\mathrm{Vir}_c$ at the
Brown--Henneaux central charge $c = 3\ell/(2G)$.  The shadow
depth is $r_{\max} = \infty$ (class~$\mathbf{M}$),
reflecting the infinite tower of soft graviton interactions.
The shadow partition function
$Z^{\mathrm{sh}} = \exp(\sum F_g\,\hbar^{2g})$
converges absolutely (Bernoulli decay), providing a
well-defined genus expansion for celestial amplitudes at all
loop orders.
\end{remark}
